% !TEX program = xelatex
% On the Design, Implementation, and Execution Environment
% of the Typed AI Agent Language AIPL

\documentclass[11pt,a4paper]{article}

\usepackage[a4paper,margin=22mm]{geometry}
\usepackage{xeCJK}
\usepackage{fontspec}
\usepackage{amsmath,amssymb,amsthm}
\usepackage{listings}
\usepackage{xcolor}
\usepackage{booktabs}
\usepackage{tabularx}
\usepackage{longtable}
\usepackage{array}
\usepackage{multirow}
\usepackage{hyperref}
\usepackage{enumitem}
\usepackage{makecell}
\usepackage{mathtools}

\setCJKmainfont{Hiragino Mincho ProN}
\setCJKsansfont{Hiragino Sans}
\setCJKmonofont{Hiragino Sans}
\setmonofont{Menlo}[Scale=0.82]

\hypersetup{
  colorlinks=true,
  linkcolor=black,
  urlcolor=blue!60!black,
  citecolor=blue!50!black,
  pdftitle={On the Design, Implementation, and Execution Environment of the Typed AI Agent Language AIPL},
  pdfauthor={Yasushi Kodama}
}

% ---- listings: AIPL language definition ----------------------------
\definecolor{codebg}{RGB}{248,248,242}
\definecolor{codecmt}{RGB}{120,120,120}
\definecolor{codekw}{RGB}{30,80,160}
\definecolor{codestr}{RGB}{160,40,40}
\definecolor{codetype}{RGB}{20,120,90}

\lstdefinelanguage{aipl}{
  keywords={class, method, function, return, var, new, send, now, future, await,
            reply, become, select, case, scope, saga, step, compensate, channel,
            if, else, while, do, self, sender, where, pub, linear, transient,
            session_define, grant_cap, revoke_cap, true, false},
  keywordstyle=\color{codekw}\bfseries,
  morekeywords=[2]{int, float, string, bool, unit, actor, array, tuple, record,
                   image, any},
  keywordstyle=[2]\color{codetype},
  comment=[l]{//},
  commentstyle=\color{codecmt}\itshape,
  string=[b]",
  stringstyle=\color{codestr},
  morestring=[b]',
  sensitive=true,
  basicstyle=\ttfamily\small,
  backgroundcolor=\color{codebg},
  frame=single,
  framesep=4pt,
  xleftmargin=4pt,
  xrightmargin=4pt,
  breaklines=true,
  showstringspaces=false,
  literate={→}{{$\rightarrow$}}1 {⊢}{{$\vdash$}}1 {∀}{{$\forall$}}1,
}
\lstset{language=aipl}

% ---- symbols -------------------------------------------------------
\newcommand{\Yes}{\textcolor{green!50!black}{\textbf{\checkmark}}}
\newcommand{\No}{\textcolor{red!70!black}{$\times$}}
\newcommand{\Part}{\textcolor{orange!80!black}{$\triangle$}}
\newcommand{\NA}{\textcolor{gray}{--}}
\newcommand{\code}[1]{\texttt{\small #1}}

\newtheorem{theorem}{Theorem}
\newtheorem{proposition}{Proposition}
\newtheorem{definition}{Definition}

\title{\bfseries On the Design, Implementation, and Execution\\
  Environment of the Typed AI Agent Language AIPL}
\author{Yasushi Kodama\\[2pt]
  {\small Hosei University\quad\texttt{yass@hosei.ac.jp}}}
\date{2026-06-12}

\begin{document}
\maketitle

\begin{abstract}
AI systems built around large language models (LLMs) are rapidly
converging on architectures in which multiple autonomous
agents---planners, solvers, reviewers, tool runners---cooperate
asynchronously by exchanging messages.
This paper presents, in a single coherent account, the design,
implementation, and execution environment of \textbf{AIPL}
(Actor-based Intelligent Parallel Language), a typed AI agent
language that expresses this structure as a first-class language
construct. AIPL inherits the lineage of Hewitt's actor model (1973)
and ABCL/1 (1986), onto which it layers (i) \textbf{Hindley--Milner
(HM) type inference} that demands no annotations, (ii) \textbf{gradual
typing features} such as capability effects, linear types and session
types, and (iii) \textbf{first-class AI integration} that embeds LLM
calls as the language primitive \code{ai\_call}.
On the implementation side, a single \code{.abcl} source is served by
\textbf{seven runtimes} (annotated/inferred Python, the canonical
OCaml, browser/Node JavaScript, and C) and \textbf{nine code-generation
backends} (C+pthread, SDL2, Xinu, Python, Pony, Erlang, Go, Prolog,
LLVM/OpenMP). We empirically verify that the three HM-equipped
implementations (OCaml, inferred Python, and the C code generator)
produce \emph{identical inferred types} for the same source.
For execution, the same source can be deployed across three concurrency
models---1:1 OS threads, M:N lightweight threads, and a cooperative
event loop---and we add a distributed runtime layer (DR-1..13) providing
token-budget governance, checkpointing, quarantine, CRDT state and saga
compensation, together with on-device distributed execution that spans
a Mac and the bare-metal OS Xinu running on a Raspberry Pi.
As evaluation we report cross-validation of type inference, a locality
analysis of type soundness, and more than 213 green smoke tests/demos.
\end{abstract}

\tableofcontents
\vspace{1em}

% =====================================================================
\section{Introduction}
% =====================================================================

\subsection{Background and Motivation}

As LLM capabilities improve, practical AI systems increasingly rely not
on a single model call but on multiple agents with distinct roles
(planner, solver, reviewer, tool executor) that \emph{cooperate
asynchronously through messages}. Strikingly, this is structurally
isomorphic to the vision of \textbf{ABCL/1}~\cite{abcl1}, the
object-oriented concurrent computation model proposed by Yonezawa et
al.\ in 1986, in which independent stateful objects coordinate through
three forms of message passing (past, now, future).
However, the actor languages of that era lacked any facility for
absorbing, at the language level, the \emph{token-budget management,
timeouts, providers, fallbacks and cross-machine coordination} that
modern AI agents require.

This work designs and realizes \textbf{AIPL} (Actor-based Intelligent
Parallel Language) as a language that fills this gap.
Building on a direct re-implementation of ABCL, AIPL aims to:
\begin{enumerate}[leftmargin=1.6em,itemsep=2pt]
  \item realize an \textbf{annotation-free actor language} via
        Hindley--Milner type inference;
  \item make LLM invocation a first-class primitive,
        \code{ai\_call(provider, prompt)}; and
  \item drive, from one source, three concurrency models
        (OS threads, lightweight threads, event loop) and beyond,
        up to distributed and on-device execution.
\end{enumerate}
AIPL is therefore not merely a processor but a ``scripting layer for an
AI-OS''---a language layer through which the runtime governs the
concurrency, persistence, budgets and cross-machine coordination of
many AI agents.

\subsection{Contributions}

This paper makes four contributions.

\begin{description}[leftmargin=1.4em,style=nextline,itemsep=3pt]
\item[(C1) An annotation-free actor language]
  We design an actor language that requires no type annotations in its
  surface syntax, inferring field types and method signatures from call
  sites and initializers, and we empirically verify that three
  independent HM implementations produce \emph{identical inference
  results} for the same source (\S\ref{sec:eval}).
\item[(C2) HM-driven code specialization]
  Using inference results, we establish a code-generation technique that
  \emph{specializes} regions outside actor message boundaries from a
  \code{value\_t} tagged union to native C types (\S\ref{sec:codegen}).
\item[(C3) First-class AI integration]
  An LLM can be driven from any of ABCL/1's three send forms, and the
  runtime governs token budgets, concurrency, fallback and usage
  accounting (\S\ref{sec:llm}).
\item[(C4) One source, many concurrency models, distributed execution]
  The same \code{.abcl} source is deployed unmodified across three
  concurrency models, and we further demonstrate on-device distributed
  execution spanning a Mac and the bare-metal OS Xinu on a Raspberry Pi
  over HTTP/JSON and UART/TCP (\S\ref{sec:runtime},\S\ref{sec:dist}).
\end{description}

\subsection{Organization}
\S\ref{sec:related} reviews background and related work,
\S\ref{sec:design} the language design, and
\S\ref{sec:type} the type system and inference.
\S\ref{sec:impl} covers the realization of seven runtimes and nine
backends; \S\ref{sec:runtime}--\S\ref{sec:dist} cover the execution
environment (concurrency models, LLM integration, distributed runtime,
evolutionary pipeline); \S\ref{sec:case} gives case studies;
\S\ref{sec:eval} evaluation, \S\ref{sec:discuss} discussion and
limitations, and \S\ref{sec:concl} concludes.

% =====================================================================
\section{Background and Related Work}\label{sec:related}
% =====================================================================

\paragraph{The actor model and ABCL.}
The actor model originates with Hewitt~\cite{hewitt}, treating each
actor as a unit of computation that (i) receives a message,
(ii) updates internal state, and (iii) creates new actors and sends.
ABCL/1~\cite{abcl1} by Yonezawa et al.\ introduced three send forms---%
\emph{past} (fire-and-forget), \emph{now} (synchronous) and
\emph{future} (asynchronous handle)---and became a classic of
object-oriented concurrent computing. AIPL inherits these three send
forms directly (\S\ref{sec:design}).

\paragraph{Hindley--Milner type inference.}
The HM type system, beginning with Hindley~\cite{hindley} and
Milner~\cite{milner}, is guaranteed by Damas--Milner's Algorithm
\textsf{W}~\cite{damas} to infer principal types \emph{without
annotations}. AIPL's inferencer adopts a textbook HM comprising
\code{unify} / \code{generalize} / \code{instantiate} and let
polymorphism (\S\ref{sec:infer}).

\paragraph{Gradual typing, effects, linearity, sessions.}
Siek--Taha's gradual typing~\cite{gradual} enables a continuous
coexistence of static and dynamic typing. While AIPL is centered on its
inferred dialect, the annotated dialect retains capability
effects~\cite{capability}, linear types~\cite{linear} and session
types~\cite{session} as a reference implementation, integrating the two
as the inferred-HM axis and the annotated-gradual axis of one language
(\S\ref{sec:gradual}).

\paragraph{Structured concurrency and distribution.}
Structured concurrency via \code{scope \{ \dots \}} shares its semantics
with Java Project Loom's \textsf{StructuredTaskScope}~\cite{loom}.
At the distributed layer, AIPL incorporates CRDT-based~\cite{crdt} actor
state replication and saga-style compensating transactions.

\paragraph{Evolutionary language design.}
AIPL's next-generation feature set was determined through a unique loop
(\S\ref{sec:evo}): an evolutionary search based on
MAP-Elites~\cite{mapelites} empirically derives axes of language design,
and the winners are implemented as features.

\paragraph{Comparison with existing actor/agent platforms.}
Erlang/OTP, Akka, Pony and Ray are mature actor/distributed platforms,
but they pursue different design goals with respect to LLM integration,
type inference and backend diversity (Table~\ref{tab:compare}).
AIPL's distinctiveness lies in simultaneously achieving (i) annotation-
free HM inference over an actor language, (ii) projection from one
frontend to nine backends, and (iii) first-class governance of LLM calls
and token budgets at the language/runtime level.

\begin{table}[h]
\centering\small
\caption{Comparison of design goals with existing platforms}\label{tab:compare}
\begin{tabularx}{\linewidth}{@{}l c c c c X@{}}
\toprule
 & 3 sends & inference & multi-backend & LLM first-class & main concurrency \\
\midrule
Erlang/OTP & \Part & \No & \No & \No & M:N (BEAM) \\
Akka (Scala) & \Part & \Part & \No & \No & M:N (JVM) \\
Pony & \No & \Part & \No & \No & M:N (ref.\ caps) \\
Ray (Python) & \Part & \No & \No & \Part & 1:1 + cluster \\
\textbf{AIPL} & \Yes & \Yes & \Yes & \Yes & 1:1 / M:N / event-loop \\
\bottomrule
\end{tabularx}
\end{table}

% =====================================================================
\section{The Design of AIPL}\label{sec:design}
% =====================================================================

\subsection{Computation Model}
An AIPL program is a set of \code{class} declarations. Each
\code{class} defines one actor species, and \code{new C(args)} creates
an instance (actor). Each actor owns an independent \textbf{mailbox},
and its fields (state) are closed within the actor and cannot be read or
written directly from outside. The sole interaction between actors is
message passing.

\subsection{The Three Send Forms}
Inheriting ABCL/1, AIPL provides three send forms
(Listing~\ref{lst:send}).

\begin{lstlisting}[caption={Three send forms and reply},label={lst:send},float=h]
class Counter {
  var count = 0;
  method tick() { count = count + 1; print("count = " + count); }
  method get()  { reply(count); }     // respond to a synchronous send
}
var c = new Counter();
send   c.tick();      // past   : fire-and-forget
var v = now c.get();  // now    : blocks until reply arrives
var f = future c.get();   // future : asynchronous handle
var w = await(f);         // obtain the value from the handle
\end{lstlisting}

\code{reply(value)} fulfills the reply slot of a now/future send.
\code{future\_done(f)} polls non-blockingly, and \code{await(f)} blocks
to obtain the value. Within a method body, the special variables
\code{self} (the actor) and \code{sender} (the sender) are available.

\subsection{become and select}
\code{become C(args)} replaces an actor's class and fields at runtime
(a representation of state transition). \code{select \{ case \dots \}}
performs selective reception from multiple channels/messages,
corresponding to Erlang's \code{receive}.

\subsection{Channels and Structured Concurrency}
\code{channel[T]} / \code{channel[T,cap=N]} provide CSP-style typed
channels (\code{channel\_send} / \code{channel\_recv} /
\code{select\_recv}). A \code{scope \{ \dots \}} block is a structured-
concurrency construct that automatically joins the futures created
within it upon exit, guaranteeing that the parent scope owns the
lifetimes of its child tasks.

\subsection{The LLM Integration Primitive}
LLM invocation is built into the language as
\code{ai\_call([provider,] prompt)}. The provider id is
\code{1}=Gemini (default) / \code{2}=Anthropic (Claude) /
\code{3}=OpenAI / \code{0}=auto, and \code{AIPL\_AI\_PROVIDER=mock}
takes all calls offline. An LLM can be driven from any of the three send
forms (Listing~\ref{lst:ai}).

\begin{lstlisting}[caption={Parallel queries via LLM $\times$ future},label={lst:ai},float=h]
class Asker { method ask(q) { var r = ai_call(2, q); reply(r); } }
var a = new Asker();
var f1 = future a.ask("Translate hello to French.");
var f2 = future a.ask("Translate hello to German.");
print(await(f1));   // the two LLM calls run in parallel
print(await(f2));
\end{lstlisting}

% =====================================================================
\section{Type System and Type Inference}\label{sec:type}
% =====================================================================

\subsection{The World of Types}
AIPL's types are as follows. Base types: \code{int} (64-bit),
\code{float} (double precision; a trailing dot is allowed, e.g.\
\code{100.}), \code{string} (UTF-8), \code{bool}, \code{unit};
structured types: \code{actor(C)} (an instance of class C),
\code{array}, \code{tuple(\dots)} (positional, immutable, fixed arity;
the length and each slot type are part of the type),
\code{record\{a:int,b:string\}} (structural, named fields),
\code{channel[T]}, function types, and \code{future}; plus implicit
polymorphism via single-letter type variables (\code{T U V \dots}).
A generic function can be written \code{function id(x:T)->T}.

\subsection{Hindley--Milner Type Inference}\label{sec:infer}
The inferencer \code{infer.ml} of the canonical implementation (OCaml)
realizes textbook HM: type variables via mutable union-find,
\code{unify} / \code{generalize} / \code{instantiate}, and let
polymorphism. The methods of a class are \textbf{pre-registered} with
fresh type variables and concretized bottom-up during body checking---a
\textbf{two-pass scheme}. As a result, information from call sites
\emph{automatically monomorphizes} method signatures.

\begin{lstlisting}[caption={No annotations = full inference; call sites fix the types},label={lst:infer},float=h]
class Hello {
  var count = 0.;        // count:float from the literal 0.
  method init(n) {       // n is unannotated
    count = n;           // count:float forces n:float
    print("hi " + n);
  }
  method greet() { print("count = " + count); }
}
var h = new Hello(5);    // re-fixes init's argument type from the call site
\end{lstlisting}

For Listing~\ref{lst:infer}, the inferencer derives
\code{count : float}, \code{init : (int) -> unit},
\code{greet : () -> unit}. A method such as \code{Counter.dec(x)} is
$\forall \alpha.\,\alpha \to \mathit{unit}$ at declaration, but is fixed
to $\alpha = \mathit{float}$ by the body \code{count - x}
(count:float), requiring no separate monomorphization pass
(HM's \code{unify} does it automatically). Note that AIPL has no
higher-rank polymorphism (rank-2 or higher); positions that take
functions as arguments fall back to \code{any}.

\subsubsection{The Two-Pass Inference Algorithm}
The central challenge in applying HM to an actor language is the
\emph{mutual reference among methods} and the \emph{sharing of fields}.
When method \code{m} of class \code{C} calls another method \code{n} of
the same class via \code{self.n(\dots)}, \code{m} must be checkable even
if \code{n}'s type is undetermined. AIPL solves this in two passes.

\begin{enumerate}[leftmargin=1.6em,itemsep=2pt]
  \item \textbf{Pass 1 (pre-registration)}: assign a fresh type variable
        to each field in the class, and a provisional signature
        $\vec{\alpha}\to\beta$ (fresh variables for arguments and
        result) to each method, registering them in the type
        environment $\Gamma$.
  \item \textbf{Pass 2 (body checking)}: check each method body
        bottom-up and solve, via \code{unify}, the equality constraints
        produced by assignments, calls and operations. Because fields
        are shared across methods, a use in one method may fix the type
        in another.
\end{enumerate}

Unification in Pass 2 is standard union-find with an occurs-check;
\code{instantiate} freshens the generalized scheme of a \code{let}
binding, and \code{generalize} closes, under $\forall$, the type
variables not free in the environment. In Listing~\ref{lst:infer},
Pass 1 gives \code{init} the type $(\alpha)\to\mathit{unit}$, and Pass 2
derives $\alpha = \mathit{float}$ from the body \code{count = n}
(with \code{count : float}). Furthermore, since the call site
\code{new Hello(5)} supplies the integer literal \code{5}, \code{init}'s
argument is finally re-fixed to \code{int}.

\subsubsection{Principal Unification Rules}
Unification $\mathsf{unify}(\tau_1,\tau_2)$ follows the rules below
(excerpt). $\sigma$ is a substitution and $\doteq$ denotes a pair being
unified.
\begin{align*}
\alpha \doteq \tau &\;\Rightarrow\; [\alpha \mapsto \tau]\quad
  (\alpha \notin \mathrm{ftv}(\tau)\ \text{occurs-check}) \\
(\tau_1 \to \tau_2) \doteq (\tau_1' \to \tau_2') &\;\Rightarrow\;
  \mathsf{unify}(\tau_1,\tau_1');\ \mathsf{unify}(\tau_2,\tau_2') \\
\mathsf{actor}(C) \doteq \mathsf{actor}(C') &\;\Rightarrow\;
  C = C' \\
\mathsf{record}(R_1) \doteq \mathsf{record}(R_2) &\;\Rightarrow\;
  \text{unify common fields}\ (\text{width subtyping}) \\
\mathsf{any} \doteq \tau &\;\Rightarrow\;
  \text{succeed}\ (\text{gradual boundary; insert transient cast})
\end{align*}
The last \code{any} rule is the junction with gradual typing and is the
locus of the soundness boundaries of \S\ref{sec:sound}.

\subsection{The Second Axis: Gradual Typing}\label{sec:gradual}
While centered on the inferred dialect, AIPL retains the annotated
dialect as a \emph{reference implementation} and provides gradual-typing
features.
\begin{itemize}[leftmargin=1.4em,itemsep=2pt]
  \item \textbf{Capability effects}:
        \code{function read\_url(u:string)->string !\{net,fs\}\{\dots\}}
        annotates method types with an effect row over
        \code{\{ai,fs,net,mut\}}. Default effects \code{BUILTIN\_EFFECTS}
        are assigned to 43 builtins and inferred.
  \item \textbf{Linear types}: \code{var fh: linear int = open("x");}
        statically guarantees single use of a resource.
  \item \textbf{Ownership}: \code{pub var holder: string = "";}
        distinguishes public fields.
  \item \textbf{Session types}:
        \code{session\_define("Solve", "solver.solve(string) ! string -> reviewer.review(\dots) -> \dots")}
        dynamically defines an inter-actor protocol.
  \item \textbf{Refinement}: predicate constraints over integers, reals
        and rationals are delegated to the SMT solver Z3 (detecting
        vacuously-false predicates at declaration time).
\end{itemize}
These exist in the inferred dialect only through the runtime, while the
annotated dialect carries the static-checking reference implementation;
this is the two-axis structure.

\subsection{Type Soundness}\label{sec:sound}
Type soundness is discussed via the two propositions of progress and
preservation.
\begin{definition}[Progress]
Execution of a well-typed program does not get stuck due to a type
mismatch (e.g.\ an arithmetic operation on a string).
\end{definition}
\begin{definition}[Preservation]
After taking an evaluation step, a value's type remains compatible with
its declared type (under a simple partial order that includes
\code{any}).
\end{definition}

Being a gradually typed language, AIPL exhibits the typical property of
being \emph{locally sound but globally unsound}. The boundaries at which
soundness may break are four: (1) inflow of \code{any} from unannotated
to annotated regions; (2) changes of class shape via dynamic reflection
(\code{compile} / \code{become} / \code{add\_method}); (3) divergence of
the recipient's signature under dynamic dispatch; and (4) rewriting an
actor field via a record. For the highest-priority unsoundness gap,
the \code{any}$\to$\code{T} boundary, we introduced a \textbf{transient
cast} via the \code{--transient} flag, inserting at the three boundaries
(variable declaration, method argument, function argument) a runtime
check governed by the same compatibility rule as the static check.

% =====================================================================
\section{Realization: Seven Runtimes and Nine Backends}\label{sec:impl}
% =====================================================================

\subsection{The Multi-Implementation Architecture}
AIPL has \textbf{seven runtime implementations} of a single source
language (Table~\ref{tab:runtimes}). Of these, HM inference is adopted by
three: OCaml, inferred Python, and the C code generator.

\begin{table}[h]
\centering\small
\caption{AIPL's seven runtime implementations}\label{tab:runtimes}
\begin{tabularx}{\linewidth}{@{}l l X@{}}
\toprule
Short name & Source & Type system \\
\midrule
Py-A (annotated) & \code{python-aipl/} & Gradual typing (Phase 11--17 reference, most complete) \\
Py-I (inferred) & \code{python-aipl-inferred/} & \textbf{HM inference} (interp shared with Py-A, 899 lines) \\
OCaml (canonical) & \code{src/*.ml} & \textbf{HM inference} (\code{infer.ml} 484 lines, WebSocket built in) \\
JS-O (via OCaml) & \code{src/app.js} + gateway & = OCaml HM (\code{/api/typecheck}) \\
JS-B (browser) & \code{browser-abcl/} & flow-sensitive light inference (\code{typecheck.js} 308 lines) \\
JS-N (Node) & \code{node-aipl-server/} & flow-sensitive (inherited), HTTP wrap \\
C (\code{aipl2c}) & \code{aipl2c.ml} + C runtimes & \textbf{HM + code specialization} \\
\bottomrule
\end{tabularx}
\end{table}

\subsection{Parsing Mechanism}
The canonical OCaml implementation performs lexing/parsing with
\code{lexer.mll} (ocamllex) + \code{parser.mly} (Menhir), feeding
\code{infer.ml} → \code{eval\_thread.ml}. The Python implementation uses
a Lark grammar (\code{grammar.lark}) and passes the AST to
\code{aipl\_typeck} / \code{aipl\_inference}. The browser/Node
implementations share browser-abcl's hand-written parser.

\subsection{HM-Driven Code Specialization}\label{sec:codegen}
The core technique of the C backend (\code{aipl2c}) is \textbf{code
specialization} using HM inference results. If \code{l\_a, l\_b} are
inferred to be \code{int}, then instead of the \code{value\_t}
tagged-union operation \code{v\_binop("+",l\_a,l\_b)} it emits the native
\code{long l\_x = (l\_a + l\_b);}. \code{value\_t} is kept uniform only at
actor message boundaries (\code{value\_t args[]}), while the actor
interior is specialized. This reconciles \emph{mailbox uniformity} with
\emph{static specialization of the interior}.

\subsection{Nine Code-Generation Backends}\label{sec:backends}
The OCaml \code{aipl2c} generates code from a single frontend to several
backends (Table~\ref{tab:backends}). A send is projected onto each
target's concurrency primitive: C+pthread emits \code{enqueue(\dots)},
Pony \code{obj.method(args)}, Erlang \code{Pid ! \{method,Args\}}, Go
\code{obj.mailbox <- msgT\{\dots\}}, and Prolog
\code{thread\_send\_message(\dots)}. Pony/Go/Erlang use a two-stage
initialization that separates construction from the init body.

\begin{table}[h]
\centering\small
\caption{Nine backends from a single frontend}\label{tab:backends}
\begin{tabularx}{\linewidth}{@{}l l X@{}}
\toprule
Flag & Target & Concurrency model \\
\midrule
(default) & C + pthread & 1:1 OS threads \\
\code{-lSDL2} & C + SDL2 GUI & 1:1 (with rendering) \\
\code{--xinu} & Xinu C & bare-metal cooperative \\
\code{--python} & Python & 1:1 OS threads \\
\code{--pony} & Pony & M:N lightweight \\
\code{--erlang} & Erlang/BEAM & M:N lightweight \\
\code{--go} & Go goroutine & M:N lightweight \\
\code{--prolog} & SWI-Prolog & threads \\
\code{--llvm / --openmp} & LLVM / OpenMP & hot/cold annotation, parallel for \\
\bottomrule
\end{tabularx}
\end{table}

% =====================================================================
\section{Execution Environment (1): Concurrency Models and LLM Integration}\label{sec:runtime}
% =====================================================================

\subsection{Unifying Three Concurrency Models}
A key demonstration in AIPL is that the same \code{.abcl} source can be
deployed, \emph{unmodified}, across three concurrency models.
\begin{enumerate}[leftmargin=1.6em,itemsep=2pt]
  \item \textbf{1:1 OS threads} (Python/OCaml/C):
        on the order of 100--1000 actors.
  \item \textbf{M:N lightweight threads} (Pony/Erlang/Go codegen):
        on the order of $10^5$--$10^7$ actors.
  \item \textbf{Cooperative event loop} (browser-abcl/node):
        single-threaded \code{setTimeout(0)} scheduling.
\end{enumerate}
Because the actor abstraction is independent of the concurrency model,
one can migrate from a prototype (event loop) to high throughput (M:N)
without changing the source.

\subsection{LLM Governance}\label{sec:llm}
The \code{ai\_call} family (\code{ai\_call\_with\_system},
\code{ai\_call\_retry}, \code{ai\_call\_image}, \code{ai\_call\_priority})
is governed by the runtime through environment variables:
\code{ABCL\_AI\_TOKEN\_BUDGET} (\code{BudgetExceeded} on overflow),
\code{ABCL\_AI\_MAX\_CONCURRENT} (in-flight semaphore),
\code{ABCL\_AI\_FALLBACK\_MODELS} (fallback chain),
\code{ABCL\_AI\_USAGE\_FILE} (usage accounting), and
\code{AIPL\_AI\_PROVIDER=mock} (offline smoke).
Backends support Ollama (local), OpenAI, Gemini and Anthropic.

\subsection{WebSocket and Dashboards}
All seven runtimes provide WebSocket over a unified wire protocol:
OCaml implements RFC 6455 by hand in \code{web\_gateway.ml} (40 lines),
Python in \code{aipl\_websocket.py}, JS-B uses the browser native API,
JS-N uses \code{ws}, and C uses libwebsockets. A client joins via
\code{?sid=<group>} and receives
\code{\{"kind":"welcome","sid":\dots,"peers":N\}}.
\code{aipl\_main.py --dashboard 8800} provides live counters, an observed
traffic table and an SSE event stream, and
\code{ABCL\_PEER\_DASHBOARDS} aggregates multiple workers into a TOTAL
row.

% =====================================================================
\section{Execution Environment (2): Distributed Runtime and On-Device Distribution}\label{sec:dist}
% =====================================================================

\subsection{Distributed Runtime DR-1..13}
Enabled by \code{AIPL\_DIST\_ENABLE=1}, the distributed runtime layer
(\code{aipl\_dist}) provides, opt-in, the governance features needed for
operating AI agents (Table~\ref{tab:dr}).

\begin{table}[h]
\centering\small
\caption{Distributed runtime features DR-1..13}\label{tab:dr}
\begin{tabularx}{\linewidth}{@{}l l X@{}}
\toprule
ID & Feature & Summary \\
\midrule
DR-1 & env-var routing & \code{AIPL\_ROUTE}, actor-name→tag table \\
DR-2 & structured log & ND-JSON, one line per event \\
DR-3 & token-budget coordination & 60-second sliding window (RPM/TPM) \\
DR-4 & checkpoint/resume & atomic save/restore \\
DR-5 & quarantine \& skip & automatic quarantine (TTL) \\
DR-6 & quorum replication & parallel multi-provider \\
DR-7 & subtree quarantine & OTP-style blast-radius limiting \\
DR-8 & spawn-tree tracking & TLS-auto parent \\
DR-9 & runtime hooks & interp/actor/call\_ai (opt-in) \\
DR-10 & CRDT state & G-Counter/OR-Set/LWW-Register (15 prims) \\
DR-11 & saga compensation & \code{saga\{step\{\}compensate\{\}\}} (LIFO) \\
DR-12 & multi-region failover & per-region route + chain walk \\
DR-13 & auto-scaling pool & hysteresis-driven (\code{pool\_create} etc.) \\
\bottomrule
\end{tabularx}
\end{table}

\subsection{HTTP/JSON Distribution and Remote Send}
Remote send is realized over an HTTP/JSON wire protocol:
\code{POST /api/json/send} (fire-and-forget) and
\code{POST /api/json/call} (synchronous reply). Python provides
\code{remote\_call} / \code{remote\_now} / \code{remote\_future} and
\code{web\_listen(port)} / \code{web\_expose} / \code{serve\_forever}.
Setting \code{ABCL\_REMOTE\_SECRET} makes an HMAC-SHA256 signature
(\code{X-ABCL-Sig}) mandatory, wire-compatible across Python/OCaml/C
(OCaml uses a pure \code{hmac\_sha256.ml}). This enables coordination
across three nodes with different providers (coordinator / solver /
verifier).

\subsection{On-Device Distribution: Mac $\times$ Raspberry Pi (Xinu)}
AIPL further demonstrates on-device distribution spanning host actors on
a Mac and AIPL actors on the bare-metal OS \textbf{Xinu} running on a
Raspberry Pi. The two are connected by a line-oriented RPC over UART/TCP
(opcodes: PING/LIST/SEND/QUERY/LOAD/COMPILE/RUN/SPAWN), and the inferred
Python translates this transparently into sends via the \code{uart1://}
scheme. As a representative example, the ``distributed dining
philosophers'' was composed of 3 Mac nodes + 2 Xinu nodes, with the 5
forks turned into Xinu actors, achieving mutual-exclusion synchronization
across the boundary on real hardware.

\subsection{The Evolutionary Language-Design Pipeline}\label{sec:evo}
AIPL's next-generation features (CE-1..13 for inference enhancement and
DR-1..13 for distribution) were implemented as the result of a
\textbf{MAP-Elites} genetic algorithm in \code{aice-evolution-v2} that,
taking a language specification (\code{.aice}) as input, performs
cross-task evaluation and LLM-directed mutation to empirically derive
``universal axes of language design.'' For example, structured
concurrency (\code{scope\{future\dots\}}) was turned into a language
feature based on a prediction in which the search pinned
\code{concurrency=structured}, \code{ownership=linear} and
\code{type=dependent} as top axes. The pipeline consists of three stages,
\code{.aice} → \code{.ga.json} → \code{.abcl}. Table~\ref{tab:evo} shows
the predicted axes the search derived and reflected into the
implementation.

\begin{table}[h]
\centering\small
\caption{Language axes derived by the evolutionary search and their implementation}\label{tab:evo}
\begin{tabularx}{\linewidth}{@{}l l X@{}}
\toprule
Prediction (axis) & Winning value & Implementation \\
\midrule
type\_safety & dependent & refinement + Z3 (CE-6/CE-9) \\
effect & capability & capability effect rows (CE-10/CE-11) \\
concurrency & csp\_channels & \code{channel[T]} (CE family / Phase 13) \\
ownership & linear & linear types \code{linear} \\
state & symbol\_owned & actor-closed fields \\
concurrency (again) & structured & \code{scope\{future\dots\}} \\
\bottomrule
\end{tabularx}
\end{table}

The winning individuals (balanced / hang-resilience / Erlang-OTP-like)
were selected from a MAP-Elites search of 8 seeds $\times$ 30 generations
$\times$ 2 runs (about 24 minutes each, over 1400 LLM calls). These
became the implementation basis for DR-1..13.

% =====================================================================
\section{Case Studies}\label{sec:case}
% =====================================================================
This section shows, through two case studies, how AIPL's design
crystallizes in real applications.

\subsection{An LLM Multi-Agent Pipeline}
A configuration in which planner, solver and reviewer agents cooperate to
solve a problem can be written concisely with the three send forms and
\code{ai\_call} (Listing~\ref{lst:pipeline}). Future-form sends overlap
the LLM calls of the solver and reviewer, and \code{scope} bundles the
lifetimes of the child tasks.

\begin{lstlisting}[caption={planner→solver→reviewer cooperation},label={lst:pipeline},float=h]
class Solver   { method solve(q)      { reply(ai_call(2, "Solve: " + q)); } }
class Reviewer { method review(q, a)  { reply(ai_call(2, "Grade:" + q + a)); } }
class Planner {
  var s = new Solver();
  var r = new Reviewer();
  method run(q) {
    scope {
      var fa = future s.solve(q);     // solver's LLM call
      var a  = await(fa);
      var fr = future r.review(q, a); // reviewer's LLM call
      reply(await(fr));
    }                                 // unfinished futures auto-joined on scope exit
  }
}
var p = new Planner();
print(now p.run("integral of x^2"));
\end{lstlisting}

The token budget (\code{ABCL\_AI\_TOKEN\_BUDGET}), concurrency
(\code{ABCL\_AI\_MAX\_CONCURRENT}), quorum replication (DR-6) and
fallback (\code{ABCL\_AI\_FALLBACK\_MODELS}) are all governed by
environment variables and the runtime without changing this source.

\subsection{Mac $\times$ Xinu On-Device Distribution: Dining Philosophers}
We ran the classic dining-philosophers problem in a mixed topology that
places the 5 forks as Xinu actors on a Raspberry Pi and the philosophers
as actors on a Mac. The inferred Python on the Mac side translates sends
to the fork actors transparently into line-oriented RPC (QUERY/SEND)
under the \code{uart1://} scheme, achieving mutual-exclusion
synchronization across the boundary on real hardware. The fact that the
same philosopher logic can be deployed \emph{unmodified} into the three
forms (a) single-process 1:1 threads, (b) M:N via Go-goroutine codegen,
and (c) Mac+Xinu on-device distribution, concretely substantiates the
concurrency-model independence of \S\ref{sec:runtime}.

% =====================================================================
\section{Evaluation}\label{sec:eval}
% =====================================================================

\subsection{Cross-Validation of Inference}
AIPL's principal empirical claim is that the three HM-equipped
implementations (OCaml / inferred Python / C code generator) produce
\emph{identical inferred types} for the same \code{.abcl} source.
A representative example follows (identical across all three):
\begin{verbatim}
Hello:   count : float;  init : (int) -> unit;
         greet : () -> unit;  inc : () -> unit
Counter: count : float;  inc : () -> unit;
         dec : (float) -> unit   <- monomorphized via call site
\end{verbatim}

\subsection{Feature Parity and Testing}
The next-generation features CE-1..13 + DR-1..13 are confirmed to have
feature parity across the main runtimes. At the completion of Phase O,
distributed DR-1..9 and inference CE-1..9 were each 9/9 across the Py-A /
OCaml / JS-O implementations, totaling 18/18 (about 1305 LOC of
implementation). JS-B / JS-N also closed out at 26/26, including
CE-10/CE-12/CE-13/DR-11. Table~\ref{tab:smoke} summarizes the smoke
tests.

\begin{table}[h]
\centering\small
\caption{Smoke-test summary (representative figures)}\label{tab:smoke}
\begin{tabularx}{\linewidth}{@{}X r@{}}
\toprule
Item & Result \\
\midrule
OCaml runtime & 52/52 \\
Browser (syntax/parse) & 7/7 + 4/4 \\
Python (\code{.abcl} / legacy \code{.aipl}) & 33/33 + 21/21 \\
Distributed cross-language (3-node mock) & 8/8 \\
Self-hosting tower (10 levels, 41 samples) & 47/47 \\
\code{aipl\_dist} unit & 27/27 \\
Distributed MVP demos (8 features $\times$ 3) & 24/24 \\
next-gen / self-host assertions & 103 all PASS \\
\midrule
\textbf{Total} & \textbf{213+ tests/demos green} \\
\bottomrule
\end{tabularx}
\end{table}

\subsection{Self-Hosting}
The \code{aipl-self-host/} tower, which describes AIPL in AIPL itself,
consists of 10 levels (Level A--C-3 plus the integrated Level Z), and all
47/47 smoke tests are green. Level Z is an integrated full self-host
using an 18-line Python bootstrap and a CE-11 capability-checked
\code{io\_bridge.abcl}, applying each level's checker to itself to
demonstrate phase monotonicity.

\subsection{Implementation Size}
The sizes of the principal components are: OCaml HM $\approx$985 lines
(\code{infer.ml} 484 + \code{types.ml} + \code{typing\_env.ml}), inferred
Python 899 lines, \code{aipl\_inference.py} $\sim$1255 lines,
\code{aipl\_dist.py} 591 lines, \code{typecheck.js} 308 lines, WebSocket
40 lines, the C+SDL2 runtime 1178 lines, and the self-host with 9 modules
$\sim$1500 lines.

% =====================================================================
\section{Discussion and Limitations}\label{sec:discuss}
% =====================================================================

\paragraph{Limits of soundness.}
As stated in \S\ref{sec:sound}, AIPL is a gradually typed language: it is
locally sound for annotated regions, but global soundness does not hold
at boundaries with unannotated regions or dynamic reflection. The
transient cast plugs the highest-priority gap, but the runtime type cast
is limited and constitutes a weak guarantee that falls short of full
transient soundness. Soundness guarantees for programs containing dynamic
reflection (\code{become} / \code{compile}) are future work.

\paragraph{Expressiveness of inference.}
The current HM lacks higher-rank polymorphism (rank-2 or higher), and
positions taking functions as arguments fall back to \code{any}. Also,
narrowing is limited to simple conditions; narrowing over \code{\&\&}
chains and nested conditions is unsupported.

\paragraph{Version drift.}
The figures in this paper are integrated from several reports of
different periods, with version drift in the number of runtimes (7 vs 6,
depending on how Py-A is counted) and the number of self-hosting levels
(9/37 vs 10/47). This paper adopts the latest feature matrix (10 levels,
47/47, 7 runtimes).

\paragraph{Future work.}
(i) Strengthening soundness guarantees for programs with dynamic
reflection; (ii) higher-rank polymorphism and full flow-sensitive
narrowing; (iii) large-scale (>$10^6$ actors) on-device evaluation on the
M:N backends; and (iv) quantitative validation of language-feature
prediction via the evolutionary pipeline.

% =====================================================================
\section{Conclusion}\label{sec:concl}
% =====================================================================
This paper presented, in a single coherent account, the design,
implementation and execution environment of AIPL, a typed AI agent
language that gives a first-class linguistic expression to the
``cooperating multiple agents'' that form the natural architecture of
LLM-driven AI systems. AIPL layers HM type inference and gradual-typing
features, together with first-class AI integration via \code{ai\_call},
onto ABCL/1's three send forms; it drives seven runtimes, nine backends
and three concurrency models from one source, and realizes on-device
distribution spanning a Mac and a Raspberry Pi (Xinu) over HTTP/JSON and
UART/TCP. We empirically showed that the three HM-equipped
implementations produce identical inferred types for the same source, and
that more than 213 tests/demos are green. AIPL offers one attainment of a
``scripting layer for an AI-OS'' that bridges the classical ideal of
actor languages and the modern operation of AI agents.

% =====================================================================
\begin{thebibliography}{99}
\bibitem{hewitt} C. Hewitt, P. Bishop, R. Steiger.
  A Universal Modular ACTOR Formalism for Artificial Intelligence.
  \textit{IJCAI}, 1973.
\bibitem{abcl1} A. Yonezawa, J.-P. Briot, E. Shibayama.
  Object-Oriented Concurrent Programming in ABCL/1.
  \textit{OOPSLA}, 1986.
\bibitem{hindley} J. R. Hindley.
  The Principal Type-Scheme of an Object in Combinatory Logic.
  \textit{Trans. AMS}, 1969.
\bibitem{milner} R. Milner.
  A Theory of Type Polymorphism in Programming.
  \textit{J. Computer and System Sciences}, 1978.
\bibitem{damas} L. Damas, R. Milner.
  Principal Type-Schemes for Functional Programs.
  \textit{POPL}, 1982.
\bibitem{gradual} J. G. Siek, W. Taha.
  Gradual Typing for Functional Languages.
  \textit{Scheme Workshop}, 2006.
\bibitem{capability} A. Mettler, D. Wagner, T. Close.
  Joe-E: A Security-Oriented Subset of Java.
  \textit{NDSS}, 2010.
\bibitem{linear} P. Wadler.
  Linear Types Can Change the World!
  \textit{Programming Concepts and Methods}, 1990.
\bibitem{session} K. Honda, V. T. Vasconcelos, M. Kubo.
  Language Primitives and Type Discipline for Structured
  Communication-Based Programming. \textit{ESOP}, 1998.
\bibitem{loom} R. Pressler et al.
  JEP 453: Structured Concurrency (Java Project Loom). 2023.
\bibitem{crdt} M. Shapiro, N. Pregui\c{c}a, C. Baquero, M. Zawirski.
  Conflict-Free Replicated Data Types. \textit{SSS}, 2011.
\bibitem{mapelites} J.-B. Mouret, J. Clune.
  Illuminating Search Spaces by Mapping Elites. arXiv:1504.04909, 2015.
\end{thebibliography}

\end{document}
